% !TEX program = xelatex
% AI-effects 프로젝트 작업 노트 --- SOC 2018–KSCO 8차 연계표 두 가지의 비교.
% 이 저장소 연계표: code/01_import_soc2018_ksco8_crosswalk.do → data/proc/exposure/bls_soc_crosswalk/
% 외부 연계표: data/raw/exposure/bls_soc_crosswalk/isco08_ksco08/SOC2018_ISCO08_KSCO8_crosswalk.xlsx
% 표 본문: results/table/22_*.tex (code/04_analysis_soc2018_ksco8_crosswalk_compare.do가 생성)
% 관련: paper 21(results/paper/21_SOC2018_KSCO8_연계표_구축.tex), report 17
% 컴파일: xelatex (kotex 필요)
\documentclass[11pt]{article}

\usepackage{fontspec}
\usepackage{kotex}
\usepackage[a4paper,margin=2.2cm]{geometry}
\usepackage{longtable}
\usepackage{booktabs}
\usepackage{array}
\usepackage{xcolor}
\usepackage{hyperref}
\usepackage{enumitem}

\hypersetup{colorlinks=true,linkcolor=blue!50!black,urlcolor=blue!50!black}
\newcolumntype{L}[1]{>{\raggedright\arraybackslash}p{#1}}
\newcolumntype{R}[1]{>{\raggedleft\arraybackslash}p{#1}}
\setlength{\LTleft}{0pt}
\setlength{\LTright}{0pt}
\renewcommand{\arraystretch}{1.15}
\setlength{\tabcolsep}{4pt}

\title{SOC 2018–KSCO 8차 연계표 비교\\[2mm]
\large --- 이 저장소 연계표 대 \texttt{SOC2018\_ISCO08\_KSCO8\_crosswalk.xlsx} ---}
\author{AI-effects 프로젝트 (작업 노트)}
\date{\today}

\begin{document}
\maketitle

\begin{abstract}
\noindent 별도 과정으로 만든 SOC 2018–KSCO 8차 연계표(\texttt{SOC2018\_ISCO08\_KSCO8\_crosswalk.xlsx}, 이하 ``외부 파일'')를 이 저장소의 연계표(\texttt{soc2018\_ksco8\_crosswalk}, paper 21)와 비교했다. 두 연계표는 같은 공식 원자료 3종을 같은 순서(SOC 2018 → SOC 2010 → ISCO-08 → KSCO 8차 세분류)로 이었다. 결과도 거의 같다. 외부 파일의 경로 2,540개는 모두 이 저장소 연계표에 있다. 3자리 ISCO 후보 시트를 포함한 수치이며, 반대 방향의 차이는 경로 5개다. 이 5개는 모두 ISCO 3314 → KSCO 2133(자연과학 시험원)을 거친다. 외부 파일의 ISCO–KSCO 시트에서 이 행의 ISCO 코드 끝에 공백이 붙어(\texttt{"3314 "}) 결합되지 않았기 때문이다. 이 때문에 외부 파일에서는 SOC 4개(15-2051, 15-2099, 19-4061, 43-9111)의 KSCO가 하나씩 적다. 그 밖의 차이는 처리 방식의 차이다. 외부 파일은 3자리 ISCO 연결을 별도 시트의 ``후보''로 두고 가중치는 만들지 않는다. 이 저장소는 3자리 연결을 주 연계표에 넣고, 평균$\cdot$배분 가중치와 KECO 2025 코드를 붙인다. 이 저장소 연계표를 그대로 쓰는 것을 권한다. 외부 파일을 쓴다면 코드의 공백을 지우고 3자리 후보 시트를 합치면 경로 집합이 이 저장소와 같아진다.
\end{abstract}

% ============================================================
\section{목적과 비교 방법}
% ============================================================

외부 파일은 사용자가 이 저장소와 다른 과정으로 만든 연계표다(\texttt{data/raw/exposure/bls\_soc\_crosswalk/isco08\_ksco08/}). 같은 목적의 연계표를 독립적으로 두 번 만든 셈이므로, 둘을 대조하면 서로를 검증할 수 있다. 본 노트는 다음을 정리한다.
\begin{enumerate}[leftmargin=*]
  \item 두 연계표의 구성과 처리 방식의 차이(2절)
  \item 단계별 연계쌍, 전체 경로, (SOC 2018, KSCO) 쌍의 일치 여부(3절)
  \item 일치하지 않는 연계의 원인과 영향(4절), SOC별 대응 수와 경로 유형 분류의 차이(5절)
\end{enumerate}
비교는 \texttt{code/04\_analysis\_soc2018\_ksco8\_crosswalk\_compare.do}로 한다. 코드는 외부 파일에 적힌 그대로 비교하고 공백을 지우지 않는다. 공백 때문에 키가 달라지는 경우를 드러내기 위해서다.

% ============================================================
\section{두 연계표의 구성}
% ============================================================

\begin{longtable}{L{3.2cm}L{6.4cm}L{6.4cm}}
\caption{두 연계표의 구성과 처리 방식}\label{tab:compose}\\
\toprule
항목 & 외부 파일 & 이 저장소 \\
\midrule
\endfirsthead
\toprule
항목 & 외부 파일 & 이 저장소 \\
\midrule
\endhead
연계 경로 & SOC 2018 → SOC 2010 → ISCO-08 → KSCO 8차 세분류(4자리) & 같음. KSCO와 1:1인 KECO 2025를 덧붙임 \\
BLS 원자료 & SOC 2010–2018 변환표, ISCO–SOC 2010 연계표(2012, 2015 갱신) & 같은 파일 \\
통계청 KSCO–ISCO 연계표 & \texttt{\_20251103085955} 판 & \texttt{\_20260127034430} 판. 세분류 시트의 (ISCO, KSCO) 쌍은 두 판이 같음(3절) \\
구성 & xlsx 시트 8개: README, 단계별 표 3개, 경로(\texttt{SOC18\_ISCO\_KSCO}), 3자리 후보, SOC 요약, 예외 & csv/dta 3종(\texttt{crosswalk}, \texttt{paths}, \texttt{stages}). do 파일로 원자료에서 다시 만듦 \\
경로 행 & 2,541행. KSCO가 없는 경로 10행(제외 ISCO 8행, 3자리 ISCO 2행)을 빈칸으로 남김 & 2,545행. KSCO까지 이어지는 경로만 남김 \\
3자리 ISCO(19-2099 → 211, 53-2022 → 315) & 주 시트에는 KSCO 없음. 별도 시트(\texttt{3digit\_ISCO\_Expanded})에 4자리 펼침 후보 9행. ``BLS가 직접 지정한 4자리 매칭이 아님''을 명시 & 같은 9쌍으로 펼쳐 주 연계표에 포함하고 \texttt{isco3\_link} = 1로 표시 \\
KSCO 표에 없는 ISCO(1113, 6310~6340) & 경로를 남기고 KSCO를 빈칸으로 둠. 예외 시트(\texttt{KSCO\_Exceptions})에 3자리 2행과 함께 10행 & 경로에서 제외(paper 21 표 12) \\
통계청 표 중복 행(2365–3119) & 그대로 둠(시트 701행) & 지움(700쌍) \\
1:N 처리 & ``가능한 연계경로''만 제공, 대표값이나 가중치 없음(README 주의 5) & 평균 가중치 \texttt{w\_mean}, \texttt{w\_flat}과 배분 가중치 \texttt{w\_alloc} \\
복수연계 표시 & BLS의 \texttt{*}(\texttt{BLS\_Part}), 통계청의 \texttt{*}(\texttt{ISCO\_KSCO\_Multi}, \texttt{KSCO\_Multi}) & 대응 수 변수(\texttt{nk\_per18}, \texttt{n18\_perk} 등)와 거친 코드 목록 \\
대분류 벗어난 연계 & \texttt{Cross\_Major\_Flag} & \texttt{ksco\_cross\_major} \\
명칭 & 원자료 그대로(BLS 각주 표시 ``(\#\#)'', 국문 ISCO 명칭 끝 공백 등이 남음) & BLS 각주 표시 삭제, 같은 ISCO의 국문 표기 통일 \\
SOC 요약 & \texttt{SOC18\_Summary}: SOC별 SOC 2010$\cdot$ISCO$\cdot$KSCO 수와 \texttt{Path\_Complexity}(1:1, 1:N, N:M) & paper 21 표 7의 경로 유형(4분류), 엑셀 \texttt{21\_soc2018\_ksco8\_crosswalk.xlsx} \\
\bottomrule
\end{longtable}

% ============================================================
\section{단계별$\cdot$경로별 일치 여부}
% ============================================================

표~\ref{tab:stages}는 같은 키로 두 연계표를 맞춘 결과다. 경로의 키는 (SOC 2018, SOC 2010, ISCO, KSCO), 쌍의 키는 (SOC 2018, KSCO)다. 외부 파일의 경로와 쌍은 KSCO가 있는 행만 센다.

\begin{longtable}{L{6.4cm}R{1.7cm}R{1.7cm}R{1.5cm}R{1.9cm}R{1.9cm}}
\caption{단계별$\cdot$경로별$\cdot$쌍별 비교}\label{tab:stages}\\
\toprule
비교 대상 & 외부 파일 & 이 저장소 & 일치 & 외부 파일에만 & 이 저장소에만 \\
\midrule
% 04_analysis_soc2018_ksco8_crosswalk_compare.do가 생성. 직접 고치지 말 것.
1. SOC 2018–SOC 2010 & 900 & 900 & 900 & 0 & 0 \\
2. SOC 2010–ISCO-08 & 1125 & 1132 & 1123 & 2 & 9 \\
3. ISCO-08–KSCO 8차 & 700 & 700 & 699 & 1 & 1 \\
\midrule
경로: SOC18\_ISCO\_KSCO 시트 & 2531 & 2545 & 2531 & 0 & 14 \\
경로: 위 + 3digit\_ISCO\_Expanded 시트 & 2540 & 2545 & 2540 & 0 & 5 \\
\midrule
(SOC 2018, KSCO) 쌍: SOC18\_ISCO\_KSCO 시트 & 2404 & 2412 & 2404 & 0 & 8 \\
(SOC 2018, KSCO) 쌍: 3자리 펼침 포함 & 2408 & 2412 & 2408 & 0 & 4 \\

\bottomrule
\end{longtable}

\begin{itemize}[leftmargin=*]
  \item \textbf{1단계}: 900쌍이 모두 같다.
  \item \textbf{2단계}: 1,123쌍이 같다. 외부 파일에만 있는 2쌍은 BLS 원자료의 3자리 연결(19-2099–211, 53-2022–315)이다. 이 저장소에만 있는 9쌍은 그것을 4자리로 펼친 것이다. 처리 방식의 차이일 뿐 내용은 같다.
  \item \textbf{3단계}: 699쌍이 같고, 양쪽에 1쌍씩 다르다. 둘 다 ISCO 3314–KSCO 2133인데, 외부 파일 쪽 코드에 끝 공백이 있어 다른 키가 된다(4.1절). 공백을 무시하면 3단계 표는 완전히 같다. 따라서 통계청 연계표 두 판(\texttt{\_20251103085955}, \texttt{\_20260127034430})의 세분류 연계쌍도 같다.
  \item \textbf{경로}: 외부 파일의 주 시트만 보면 이 저장소에만 있는 경로가 14개다. 3자리 펼침 9개와 ISCO 3314 → KSCO 2133 5개다. 3자리 후보 시트를 합치면 차이는 5개로 줄고, 외부 파일에만 있는 경로는 없다.
  \item \textbf{(SOC 2018, KSCO) 쌍}: 3자리 후보까지 합치면 외부 파일 2,408쌍은 모두 이 저장소 2,412쌍에 들어 있다. 남는 4쌍이 ISCO 3314 경로다.
\end{itemize}

% ============================================================
\section{일치하지 않는 연계}
% ============================================================

\begin{longtable}{L{1.8cm}L{5.8cm}L{4.8cm}L{3.0cm}}
\caption{한쪽에만 있는 연계(3자리 처리 차이 제외)}\label{tab:diff}\\
\toprule
수준 & 왼쪽(경로는 SOC 2018 → SOC 2010 → ISCO) & KSCO & 있는 쪽 \\
\midrule
% 04_analysis_soc2018_ksco8_crosswalk_compare.do가 생성. 직접 고치지 말 것.
3단계 쌍 & 3314 + 끝 공백 1자 & 2133 자연과학 시험원 & 외부 파일에만 \\
3단계 쌍 & 3314 & 2133 자연과학 시험원 & 이 저장소에만 \\
경로 & 15-2051 → 15-2099 → 3314 & 2133 자연과학 시험원 & 이 저장소에만 \\
경로 & 15-2099 → 15-2091 → 3314 & 2133 자연과학 시험원 & 이 저장소에만 \\
경로 & 15-2099 → 15-2099 → 3314 & 2133 자연과학 시험원 & 이 저장소에만 \\
경로 & 19-4061 → 19-4061 → 3314 & 2133 자연과학 시험원 & 이 저장소에만 \\
경로 & 43-9111 → 43-9111 → 3314 & 2133 자연과학 시험원 & 이 저장소에만 \\
쌍 & 15-2051 & 2133 자연과학 시험원 & 이 저장소에만 \\
쌍 & 15-2099 & 2133 자연과학 시험원 & 이 저장소에만 \\
쌍 & 19-4061 & 2133 자연과학 시험원 & 이 저장소에만 \\
쌍 & 43-9111 & 2133 자연과학 시험원 & 이 저장소에만 \\

\bottomrule
\end{longtable}

\subsection{ISCO 3314 → KSCO 2133 누락(외부 파일)}

\begin{itemize}[leftmargin=*]
  \item \textbf{원인}: 외부 파일 \texttt{ISCO08\_KSCO8} 시트의 결합용 열 \texttt{ISCO08\_Code}가 이 행에서 \texttt{"3314 "}(5자, 마지막 바이트 32 = 공백)다. 같은 행의 \texttt{국제코드}는 \texttt{3314}로 깨끗하다. 반면 \texttt{국제코드2}는 \texttt{"*3314 "}다. 따라서 \texttt{국제코드2}에서 복수연계 표시 \texttt{*}만 떼고 공백은 지우지 않은 것으로 보인다(작업자 추정). 이 저장소가 쓰는 \texttt{\_20260127034430} 판의 \texttt{국제코드} 열에는 4자리가 아닌 코드가 없다.
  \item \textbf{결과}: 외부 파일 경로 시트에서 ISCO 3314(Statistical, mathematical and related associate professionals)는 KSCO 3320(통계$\cdot$데이터 관련 사무원)에만 이어진다. 통계청 표에 있는 KSCO 2133(자연과학 시험원) 연결이 빠진다.
  \item \textbf{영향받는 SOC}: Data Scientists(15-2051), Mathematical Science Occupations, All Other(15-2099), Social Science Research Assistants(19-4061), Statistical Assistants(43-9111) 4개다. 외부 파일에서는 KSCO가 하나씩 적다. 이 저장소에서는 각 SOC의 양이 KSCO 3320과 2133에 절반씩 배분된다(\texttt{w\_alloc} = 0.5). KSCO 2133 쪽에서 보면 외부 파일은 SOC 4개를 덜 받는다.
  \item \textbf{어느 쪽이 맞는가}: 통계청 공식 연계표에 3314–2133 연계가 있으므로 이 저장소 쪽이 원자료와 일치한다. 다만 통계 보조원이 자연과학 시험원과 이어지는 것이 의미상 자연스러운지는 따로 판단할 문제다. 공식 표를 따른다는 원칙상 이 저장소는 그대로 둔다.
\end{itemize}

\subsection{처리 방식만 다른 경우}

표~\ref{tab:exc}는 외부 파일의 예외 시트 10행과 이 저장소의 처리다.

\begin{longtable}{L{2.8cm}L{6.6cm}L{4.4cm}R{1.8cm}}
\caption{외부 파일 \texttt{KSCO\_Exceptions}와 이 저장소의 처리}\label{tab:exc}\\
\toprule
SOC 2018 (SOC 2010) & BLS ISCO & 이 저장소 처리 & 이 저장소 KSCO 수 \\
\midrule
% 04_analysis_soc2018_ksco8_crosswalk_compare.do가 생성. 직접 고치지 말 것.
11-1011 & 1113 Traditional chiefs and heads of villages & 경로에서 제외 & 2 \\
11-1031 & 1113 Traditional chiefs and heads of villages & 경로에서 제외 & 2 \\
19-2099 & 211 Physical and earth science professionals & 4자리로 펼쳐 연결(isco3\_link = 1) & 1 \\
53-2022 & 315 Ship and aircraft controllers and technicians & 4자리로 펼쳐 연결(isco3\_link = 1) & 3 \\
45-2092 & 6310 Subsistence crop farmers & 경로에서 제외 & 2 \\
45-2093 & 6320 Subsistence livestock farmers & 경로에서 제외 & 2 \\
45-2092 & 6330 Subsistence mixed crop and livestock farmers & 경로에서 제외 & 2 \\
45-2093 & 6330 Subsistence mixed crop and livestock farmers & 경로에서 제외 & 2 \\
45-3031 (45-3011) & 6340 Subsistence fishers, hunters, trappers and gatherers & 경로에서 제외 & 3 \\
45-3031 (45-3021) & 6340 Subsistence fishers, hunters, trappers and gatherers & 경로에서 제외 & 3 \\

\bottomrule
\end{longtable}

\begin{itemize}[leftmargin=*]
  \item \textbf{3자리 ISCO 2건}: 두 연계표 모두 같은 소분류의 ISCO 4자리 전부로 펼친다. 외부 파일은 ``KSCO 표에 나타나는 하위 4자리'', 이 저장소는 ``BLS 표에 나타나는 하위 4자리''로 펼친다. 결과는 같은 9쌍이다. 다른 점은 외부 파일이 이를 후보로 분리해 둔다는 것이다. 외부 파일의 주 시트만 쓰면 두 SOC는 KSCO가 없다.
  \item \textbf{KSCO 표에 없는 ISCO 8행}: 외부 파일은 경로를 남기고 KSCO를 빈칸으로, 이 저장소는 경로를 버린다. 해당 SOC는 모두 다른 ISCO로 KSCO와 이어진다(표의 마지막 열). 그래서 쓸 수 있는 연계는 두 연계표가 같다.
\end{itemize}

% ============================================================
\section{SOC별 대응 수와 경로 유형}
% ============================================================

외부 파일의 \texttt{SOC18\_Summary} 시트는 SOC별 SOC 2010$\cdot$ISCO$\cdot$KSCO 수를 준다. 표~\ref{tab:soccount}는 이를 이 저장소 경로에서 센 값과 비교한 것이다. 이 저장소의 ISCO 수는 KSCO까지 이어지는 ISCO만 센다.

\begin{longtable}{L{3.2cm}R{2.0cm}R{2.6cm}R{2.6cm}R{2.4cm}R{2.4cm}}
\caption{SOC 2018 867개의 대응 수 비교}\label{tab:soccount}\\
\toprule
 & 같음 & 외부 파일이 많음 & 이 저장소가 많음 & 평균(외부) & 평균(이 저장소) \\
\midrule
% 04_analysis_soc2018_ksco8_crosswalk_compare.do가 생성. 직접 고치지 말 것.
SOC 2010 수 & 867 & 0 & 0 & 1.04 & 1.04 \\
ISCO-08 수 & 860 & 5 & 2 & 1.48 & 1.48 \\
KSCO 8차 수 & 861 & 0 & 6 & 2.77 & 2.78 \\

\bottomrule
\end{longtable}

\begin{itemize}[leftmargin=*]
  \item \textbf{SOC 2010 수}: 867개 모두 같다.
  \item \textbf{ISCO 수}: 7개가 다르다. 외부 파일이 많은 5개(11-1011, 11-1031, 45-2092, 45-2093, 45-3031)는 KSCO 표에 없는 ISCO까지 세기 때문이다. 이 저장소가 많은 2개(19-2099, 53-2022)는 3자리 ISCO를 4자리로 펼쳐 세기 때문이다.
  \item \textbf{KSCO 수}: 6개가 다르고, 모두 이 저장소가 많다(표~\ref{tab:socdiff}).
\end{itemize}

\begin{longtable}{L{1.5cm}L{6.0cm}R{1.4cm}R{1.6cm}L{5.6cm}}
\caption{KSCO 수가 다른 SOC 2018}\label{tab:socdiff}\\
\toprule
SOC 2018 & 직업 & 외부 파일 & 이 저장소 & 원인 \\
\midrule
% 04_analysis_soc2018_ksco8_crosswalk_compare.do가 생성. 직접 고치지 말 것.
15-2051 & Data Scientists & 1 & 2 & ISCO 3314–KSCO 2133 연계 누락(외부 파일) \\
15-2099 & Mathematical Science Occupations, All Other & 1 & 2 & ISCO 3314–KSCO 2133 연계 누락(외부 파일) \\
19-2099 & Physical Scientists, All Other & 0 & 1 & 3자리 ISCO: 외부 파일 주 시트에 KSCO 없음(별도 시트에 후보) \\
19-4061 & Social Science Research Assistants & 1 & 2 & ISCO 3314–KSCO 2133 연계 누락(외부 파일) \\
43-9111 & Statistical Assistants & 1 & 2 & ISCO 3314–KSCO 2133 연계 누락(외부 파일) \\
53-2022 & Airfield Operations Specialists & 0 & 3 & 3자리 ISCO: 외부 파일 주 시트에 KSCO 없음(별도 시트에 후보) \\

\bottomrule
\end{longtable}

외부 파일의 \texttt{Path\_Complexity}는 SOC 쪽 경로를 1:1(347개), 1:N(489개), N:M(31개)으로 나눈다. 정의는 적혀 있지 않다. 다음과 같이 읽으면 867개 모두와 맞는다(작업자 추정, do 파일에서 assert로 확인). N:M은 SOC 2010이 둘 이상인 경우다. 1:1은 SOC 2010과 ISCO가 하나이고 KSCO가 하나 이하인 경우다. 나머지가 1:N이다. 표~\ref{tab:complex}는 이를 paper 21의 경로 유형(paper 21 표 7)과 교차한 것이다.

\begin{longtable}{L{7.0cm}R{1.8cm}R{1.8cm}R{1.8cm}R{1.8cm}}
\caption{외부 파일 \texttt{Path\_Complexity} × 이 저장소 경로 유형(SOC 수)}\label{tab:complex}\\
\toprule
이 저장소 경로 유형 & 1:1 & 1:N & N:M & 합계 \\
\midrule
% 04_analysis_soc2018_ksco8_crosswalk_compare.do가 생성. 직접 고치지 말 것.
SOC 2010 1개 → ISCO 1개 → KSCO 1개 & 342 & 0 & 0 & 342 \\
ISCO 1개 → KSCO 여러 개 & 3 & 321 & 0 & 324 \\
SOC 2010 1개 → ISCO 여러 개 & 2 & 168 & 0 & 170 \\
SOC 2010 여러 개 & 0 & 0 & 31 & 31 \\
\midrule
합계 & 347 & 489 & 31 & 867 \\

\bottomrule
\end{longtable}

\begin{itemize}[leftmargin=*]
  \item 이 저장소에서 경로 전체가 하나인 SOC 342개는 외부 파일에서도 모두 1:1이다. N:M 31개도 이 저장소의 ``SOC 2010 여러 개'' 31개와 같다.
  \item 외부 파일의 1:1이 5개 더 많다. 3개(15-2051, 19-4061, 43-9111)는 ISCO 3314 누락으로 KSCO가 하나로 줄어든 SOC다. 2개(19-2099, 53-2022)는 주 시트에 KSCO가 없어(0개) 1:1로 분류된 3자리 ISCO SOC다. 15-2099도 3314 누락의 영향을 받지만, SOC 2010이 둘이라 N:M으로 분류되어 이 차이에 들어가지 않는다.
  \item \texttt{Path\_Complexity}는 SOC 쪽 나뉨만 본다. KSCO 쪽에서도 상대가 하나인지는 보지 않는다. 양쪽 모두 1:1인 (SOC, KSCO) 쌍은 이 저장소 기준 29개뿐이다(paper 21 표 3). 외부 파일의 ``1:1 path 347''을 값을 그대로 옮길 수 있는 직업 수로 읽으면 안 된다.
\end{itemize}

% ============================================================
\section{종합과 권고}
% ============================================================

\begin{enumerate}[leftmargin=*]
  \item \textbf{독립적으로 만든 두 연계표가 사실상 같다.} 원자료, 연계 순서, 3자리 ISCO 펼침, 제외 ISCO의 결과가 모두 일치한다. 차이는 공백 하나 때문에 생긴 경로 5개뿐이다. 이 비교는 paper 21 연계표의 검증으로도 볼 수 있다.
  \item \textbf{분석에는 이 저장소 연계표를 쓴다.} 공식 표의 ISCO 3314–KSCO 2133 연계가 들어 있고, 가중치와 KECO 코드가 있다. do 파일로 원자료에서 다시 만들 수 있다.
  \item \textbf{외부 파일을 쓸 경우}: (a) \texttt{ISCO08\_Code}의 앞뒤 공백을 지우고 경로를 다시 결합한다. (b) \texttt{3digit\_ISCO\_Expanded} 시트를 주 시트에 합친다. (c) KSCO가 빈 10행을 뺀다. 이 세 단계를 거치면 경로 집합이 이 저장소와 같아진다. 가중치는 따로 정해야 한다(README 권장 사항과 같음).
  \item \textbf{외부 파일에 알려 둘 점}: 3314 행의 공백, 그리고 \texttt{Path\_Complexity}의 정의(특히 KSCO가 0개인 SOC를 1:1로 분류하는 점)를 README에 적어 두는 것이 좋다. 외부 파일은 원자료 폴더(\texttt{data/raw/})에 있으므로 이 저장소에서는 고치지 않았다.
\end{enumerate}

% ============================================================
\section{재현}
% ============================================================

Stata 17에서 \texttt{code/01\_import\_soc2018\_ksco8\_crosswalk.do}를 먼저 실행한 뒤 \texttt{code/04\_analysis\_soc2018\_ksco8\_crosswalk\_compare.do}를 실행한다. 두 번째 do 파일은 외부 파일의 시트 7개를 직접 읽어 본 노트의 표(\texttt{results/table/22\_*.tex})를 만든다.

\end{document}
